% =============================================================
% TRACEABILITY TABLE v2
% Same structure as Table 1.1 (tab:rq_obj_contrib_mapping) with
% an added "Key Evidence" column
% =============================================================
% Include via % =============================================================
% TRACEABILITY TABLE v2
% Same structure as Table 1.1 (tab:rq_obj_contrib_mapping) with
% an added "Key Evidence" column
% =============================================================
% Include via % =============================================================
% TRACEABILITY TABLE v2
% Same structure as Table 1.1 (tab:rq_obj_contrib_mapping) with
% an added "Key Evidence" column
% =============================================================
% Include via % =============================================================
% TRACEABILITY TABLE v2
% Same structure as Table 1.1 (tab:rq_obj_contrib_mapping) with
% an added "Key Evidence" column
% =============================================================
% Include via \input{traceability_table_v2.tex}
% Requires: booktabs, tabularx, lscape (or pdflscape)

\begin{landscape}
\begin{table}[htbp]
\centering
\caption{Mapping of Research Questions to Objectives, Contributions, and Key Evidence}
\label{tab:rq_obj_contrib_evidence}

\small
\setlength{\tabcolsep}{4pt}
\renewcommand{\arraystretch}{1.25}

\begin{tabularx}{\linewidth}{p{3.0cm} X X p{2.8cm} X}
\toprule
\textbf{Research Question (RQ)} &
\textbf{Primary objective(s) that directly answer it} &
\textbf{Supporting objective(s) (enablers / infrastructure)} &
\textbf{Contribution(s) most directly produced} &
\textbf{Key evidence} \\
\midrule

\textbf{RQ1:} How robust is MCQ-based evaluation to variations in distractor/option structure in systems engineering, and what risks does this pose when MCQ accuracy is used as primary qualification evidence? &
(5) Quantify MCQ robustness/sensitivity via statistical tests and effect sizes &
(1) Curate and version benchmark artifacts, including controlled MCQ variants suitable for robustness studies\newline
(3) Execute standardized inference across variants\newline
(7) Synthesize guidance framing robustness implications &
C2 (Robustness analysis of MCQ evaluation using distractor/option variation) &
Positional bias heatmaps; per-position accuracy variance; $\chi^2$ tests of answer-position dependence; position-rotated MCQ runs across all models \\
\midrule

\textbf{RQ2:} How do MCQ and OSQ modalities differ in the type and strength of evidence they provide for qualifying LLMs in systems engineering, and how can evidence from multiple modalities be combined to support task-aligned, risk-aware qualification decisions? &
(5) Quantify cross-modality evidence patterns and uncertainty\newline
(7) Synthesize qualification guidance (task-/risk-aligned fusion) &
(1) Curate and version benchmark artifacts, including paired MCQ/OSQ items suitable for cross-modality comparison\newline
(2) Generate OSQ grading scaffolds\newline
(3) Execute standardized inference\newline
(4) Score OSQ via rubric-based judging and judge stability checks &
C1 (Multi-format SysEngBench extension)\newline
C3 (Evaluation modality fusion as qualification basis) &
MCQ$\rightarrow$OSQ conversion pipeline with suitability scoring; rubric design and canonical answers; inter-judge agreement analysis; MCQ--OSQ scatter and correlation; per-category $\Delta$ heatmaps; modality classification criteria and geometric regions \\
\midrule

\textbf{RQ3:} How do evaluation modality and response length influence the computational cost of qualification, and what trade-offs arise between judged quality and token efficiency? &
(6) Evaluate token-efficiency and cost trade-offs; extend cost modeling to OSQ evaluation &
(3) Standardized inference with token/length logging\newline
(4) Rubric-based judging to produce quality signals &
C4 (Token-efficient qualification: cost-aware response length trade-offs) &
Token accounting across MCQ and OSQ modalities; response-length distributions; quality--token Pareto frontier; cost modeling under token-constrained regimes \\
\bottomrule
\end{tabularx}
\end{table}
\end{landscape}

% Requires: booktabs, tabularx, lscape (or pdflscape)

\begin{landscape}
\begin{table}[htbp]
\centering
\caption{Mapping of Research Questions to Objectives, Contributions, and Key Evidence}
\label{tab:rq_obj_contrib_evidence}

\small
\setlength{\tabcolsep}{4pt}
\renewcommand{\arraystretch}{1.25}

\begin{tabularx}{\linewidth}{p{3.0cm} X X p{2.8cm} X}
\toprule
\textbf{Research Question (RQ)} &
\textbf{Primary objective(s) that directly answer it} &
\textbf{Supporting objective(s) (enablers / infrastructure)} &
\textbf{Contribution(s) most directly produced} &
\textbf{Key evidence} \\
\midrule

\textbf{RQ1:} How robust is MCQ-based evaluation to variations in distractor/option structure in systems engineering, and what risks does this pose when MCQ accuracy is used as primary qualification evidence? &
(5) Quantify MCQ robustness/sensitivity via statistical tests and effect sizes &
(1) Curate and version benchmark artifacts, including controlled MCQ variants suitable for robustness studies\newline
(3) Execute standardized inference across variants\newline
(7) Synthesize guidance framing robustness implications &
C2 (Robustness analysis of MCQ evaluation using distractor/option variation) &
Positional bias heatmaps; per-position accuracy variance; $\chi^2$ tests of answer-position dependence; position-rotated MCQ runs across all models \\
\midrule

\textbf{RQ2:} How do MCQ and OSQ modalities differ in the type and strength of evidence they provide for qualifying LLMs in systems engineering, and how can evidence from multiple modalities be combined to support task-aligned, risk-aware qualification decisions? &
(5) Quantify cross-modality evidence patterns and uncertainty\newline
(7) Synthesize qualification guidance (task-/risk-aligned fusion) &
(1) Curate and version benchmark artifacts, including paired MCQ/OSQ items suitable for cross-modality comparison\newline
(2) Generate OSQ grading scaffolds\newline
(3) Execute standardized inference\newline
(4) Score OSQ via rubric-based judging and judge stability checks &
C1 (Multi-format SysEngBench extension)\newline
C3 (Evaluation modality fusion as qualification basis) &
MCQ$\rightarrow$OSQ conversion pipeline with suitability scoring; rubric design and canonical answers; inter-judge agreement analysis; MCQ--OSQ scatter and correlation; per-category $\Delta$ heatmaps; modality classification criteria and geometric regions \\
\midrule

\textbf{RQ3:} How do evaluation modality and response length influence the computational cost of qualification, and what trade-offs arise between judged quality and token efficiency? &
(6) Evaluate token-efficiency and cost trade-offs; extend cost modeling to OSQ evaluation &
(3) Standardized inference with token/length logging\newline
(4) Rubric-based judging to produce quality signals &
C4 (Token-efficient qualification: cost-aware response length trade-offs) &
Token accounting across MCQ and OSQ modalities; response-length distributions; quality--token Pareto frontier; cost modeling under token-constrained regimes \\
\bottomrule
\end{tabularx}
\end{table}
\end{landscape}

% Requires: booktabs, tabularx, lscape (or pdflscape)

\begin{landscape}
\begin{table}[htbp]
\centering
\caption{Mapping of Research Questions to Objectives, Contributions, and Key Evidence}
\label{tab:rq_obj_contrib_evidence}

\small
\setlength{\tabcolsep}{4pt}
\renewcommand{\arraystretch}{1.25}

\begin{tabularx}{\linewidth}{p{3.0cm} X X p{2.8cm} X}
\toprule
\textbf{Research Question (RQ)} &
\textbf{Primary objective(s) that directly answer it} &
\textbf{Supporting objective(s) (enablers / infrastructure)} &
\textbf{Contribution(s) most directly produced} &
\textbf{Key evidence} \\
\midrule

\textbf{RQ1:} How robust is MCQ-based evaluation to variations in distractor/option structure in systems engineering, and what risks does this pose when MCQ accuracy is used as primary qualification evidence? &
(5) Quantify MCQ robustness/sensitivity via statistical tests and effect sizes &
(1) Curate and version benchmark artifacts, including controlled MCQ variants suitable for robustness studies\newline
(3) Execute standardized inference across variants\newline
(7) Synthesize guidance framing robustness implications &
C2 (Robustness analysis of MCQ evaluation using distractor/option variation) &
Positional bias heatmaps; per-position accuracy variance; $\chi^2$ tests of answer-position dependence; position-rotated MCQ runs across all models \\
\midrule

\textbf{RQ2:} How do MCQ and OSQ modalities differ in the type and strength of evidence they provide for qualifying LLMs in systems engineering, and how can evidence from multiple modalities be combined to support task-aligned, risk-aware qualification decisions? &
(5) Quantify cross-modality evidence patterns and uncertainty\newline
(7) Synthesize qualification guidance (task-/risk-aligned fusion) &
(1) Curate and version benchmark artifacts, including paired MCQ/OSQ items suitable for cross-modality comparison\newline
(2) Generate OSQ grading scaffolds\newline
(3) Execute standardized inference\newline
(4) Score OSQ via rubric-based judging and judge stability checks &
C1 (Multi-format SysEngBench extension)\newline
C3 (Evaluation modality fusion as qualification basis) &
MCQ$\rightarrow$OSQ conversion pipeline with suitability scoring; rubric design and canonical answers; inter-judge agreement analysis; MCQ--OSQ scatter and correlation; per-category $\Delta$ heatmaps; modality classification criteria and geometric regions \\
\midrule

\textbf{RQ3:} How do evaluation modality and response length influence the computational cost of qualification, and what trade-offs arise between judged quality and token efficiency? &
(6) Evaluate token-efficiency and cost trade-offs; extend cost modeling to OSQ evaluation &
(3) Standardized inference with token/length logging\newline
(4) Rubric-based judging to produce quality signals &
C4 (Token-efficient qualification: cost-aware response length trade-offs) &
Token accounting across MCQ and OSQ modalities; response-length distributions; quality--token Pareto frontier; cost modeling under token-constrained regimes \\
\bottomrule
\end{tabularx}
\end{table}
\end{landscape}

% Requires: booktabs, tabularx, lscape (or pdflscape)

\begin{landscape}
\begin{table}[htbp]
\centering
\caption{Mapping of Research Questions to Objectives, Contributions, and Key Evidence}
\label{tab:rq_obj_contrib_evidence}

\small
\setlength{\tabcolsep}{4pt}
\renewcommand{\arraystretch}{1.25}

\begin{tabularx}{\linewidth}{p{3.0cm} X X p{2.8cm} X}
\toprule
\textbf{Research Question (RQ)} &
\textbf{Primary objective(s) that directly answer it} &
\textbf{Supporting objective(s) (enablers / infrastructure)} &
\textbf{Contribution(s) most directly produced} &
\textbf{Key evidence} \\
\midrule

\textbf{RQ1:} How robust is MCQ-based evaluation to variations in distractor/option structure in systems engineering, and what risks does this pose when MCQ accuracy is used as primary qualification evidence? &
(5) Quantify MCQ robustness/sensitivity via statistical tests and effect sizes &
(1) Curate and version benchmark artifacts, including controlled MCQ variants suitable for robustness studies\newline
(3) Execute standardized inference across variants\newline
(7) Synthesize guidance framing robustness implications &
C2 (Robustness analysis of MCQ evaluation using distractor/option variation) &
Positional bias heatmaps; per-position accuracy variance; $\chi^2$ tests of answer-position dependence; position-rotated MCQ runs across all models \\
\midrule

\textbf{RQ2:} How do MCQ and OSQ modalities differ in the type and strength of evidence they provide for qualifying LLMs in systems engineering, and how can evidence from multiple modalities be combined to support task-aligned, risk-aware qualification decisions? &
(5) Quantify cross-modality evidence patterns and uncertainty\newline
(7) Synthesize qualification guidance (task-/risk-aligned fusion) &
(1) Curate and version benchmark artifacts, including paired MCQ/OSQ items suitable for cross-modality comparison\newline
(2) Generate OSQ grading scaffolds\newline
(3) Execute standardized inference\newline
(4) Score OSQ via rubric-based judging and judge stability checks &
C1 (Multi-format SysEngBench extension)\newline
C3 (Evaluation modality fusion as qualification basis) &
MCQ$\rightarrow$OSQ conversion pipeline with suitability scoring; rubric design and canonical answers; inter-judge agreement analysis; MCQ--OSQ scatter and correlation; per-category $\Delta$ heatmaps; modality classification criteria and geometric regions \\
\midrule

\textbf{RQ3:} How do evaluation modality and response length influence the computational cost of qualification, and what trade-offs arise between judged quality and token efficiency? &
(6) Evaluate token-efficiency and cost trade-offs; extend cost modeling to OSQ evaluation &
(3) Standardized inference with token/length logging\newline
(4) Rubric-based judging to produce quality signals &
C4 (Token-efficient qualification: cost-aware response length trade-offs) &
Token accounting across MCQ and OSQ modalities; response-length distributions; quality--token Pareto frontier; cost modeling under token-constrained regimes \\
\bottomrule
\end{tabularx}
\end{table}
\end{landscape}
